%%%%%%%%%%%%%%%%%%%%%%%%%%%%%%%%%%%%%%%%%%%%%%%%%%%%%%%%%%%%%%%%%%%%%%%%%%%
%
% Generic template for TFC/TFM/TFG/Tesis
%
% $Id: preamble.tex,v 1.10 2013/11/05 00:23:36 macias Exp $
%
% By:
%  + Javier Mac�as-Guarasa. 
%    Departamento de Electr�nica
%    Universidad de Alcal�
%  + Roberto Barra-Chicote. 
%    Departamento de Ingenier�a Electr�nica
%    Universidad Polit�cnica de Madrid   
% 
% Based on original sources by Roberto Barra, Manuel Oca�a, Jes�s Nuevo,
% Pedro Revenga, Fernando Herr�nz and Noelia Hern�ndez. Thanks a lot to
% all of them, and to the many anonymous contributors found (thanks to
% google) that provided help in setting all this up.
%
% If you think you can add pieces of relevant/useful examples,
% improvements, please contact us at (macias@depeca.uah.es)
%
% (copyleft)2013
%
%%%%%%%%%%%%%%%%%%%%%%%%%%%%%%%%%%%%%%%%%%%%%%%%%%%%%%%%%%%%%%%%%%%%%%%%%%%

\documentclass[svgnames,spanish,openright]{book}
%\documentclass[english,openright]{book}
%\documentclass[11pt,english,twoside,openright]{book}

%\usepackage[a4,cam,center]{crop}
%\crop[font=\upshape\mdseries\small\textsf]

\usepackage[latin1]{inputenc} % Para poder escribir con acentos y �.
\usepackage[T1]{fontenc}      % Para que haga bien la ``hyphenation''. No
                              % usar si no es necesario, porque ralentiza muchisimo la compilaci�n.
\usepackage{ae}               % Para que todas las fuentes sean Type1, y ninguna Type3.

\usepackage[spanish, english]{babel}

% Use this if you want to include pdf files in the final document
\usepackage[final]{pdfpages}

% Use this if you want to delete headers and footers in empty pages
\usepackage{emptypage}

%\usepackage[nottoc]{tocbibind}
\usepackage{tocbibind}

\usepackage{listings}
\usepackage{longtable}
\usepackage{afterpage}

\usepackage{xspace}
\usepackage{verbatim}
\usepackage{moreverb}
\usepackage{multicol}
\usepackage{amsmath}
\usepackage{eurosym}
\usepackage{subfigure}
\usepackage{multirow}
\usepackage{fancyhdr}
\usepackage{makeidx}
\usepackage{rotating}
\usepackage{supertabular}
\usepackage{hhline}
\usepackage{array}
\usepackage[noadjust]{cite}      % Written by Donald Arseneau
                        % V1.6 and later of IEEEtran pre-defines the format
                        % of the cite.sty package \cite{} output to follow
                        % that of IEEE. Loading the cite package will
                        % result in citation numbers being automatically
                        % sorted and properly "ranged". i.e.,
                        % [1], [9], [2], [7], [5], [6]
                        % (without using cite.sty)
                        % will become:
                        % [1], [2], [5]--[7], [9] (using cite.sty)
                        % cite.sty's \cite will automatically add leading
                        % space, if needed. Use cite.sty's noadjust option
                        % (cite.sty V3.8 and later) if you want to turn this
                        % off. cite.sty is already installed on most LaTeX
                        % systems. The latest version can be obtained at:
                        % http://www.ctan.org/tex-archive/macros/latex/contrib/supported/cite/


\usepackage[RGB,rgb]{xcolor}
%\usepackage{color}

% To draw rectagles in tfm cover
\usepackage{tikz}


%\usepackage[authoryear]{natbib}
% \makeatletter
% \let\NAT@parse\undefined
% \makeatother
% \usepackage{natbib}

\usepackage{geometry}
\geometry{verbose,a4paper,tmargin=2.5cm,bmargin=2.5cm,lmargin=2.5cm,rmargin=2.5cm}
%\geometry{paperwidth=210mm,paperheight=297mm}

\usepackage[
ps2pdf,                %%% hyper-references for ps2pdf
bookmarks=true,%                   %%% generate bookmarks ...
bookmarksnumbered=true,            %%% ... with numbers
hypertexnames=false,               %%% needed for correct links to
                                %%% figures!!!
%hypertexnames=true,               %%% needed for correct links on pagebackrefs!!!
breaklinks=true,                   %%% breaks lines, but links are very small
%pagebackref=true,
%linktocpage=true,                 %%% enlace en el numero de p�gina.
linktoc=all,
colorlinks=true,
linkcolor=azul,    
citecolor=verde,
urlcolor=azul,                     %%% texto  con color
linkbordercolor={0 0 1},           %%% blue frames around links
pdfborder={0 0 112.0}              %%% border-width of frames 
]{hyperref}                        %%% will be multiplied with 0.009 by ps2pdf


% Para numerar las \subsubsection
\setcounter{secnumdepth}{5}
%para hacer que las \subsubsection aparezcan en el indice
\setcounter{tocdepth}{5}
%\setcounter{lofdepth}{2}
\setcounter{table}{1}
\setcounter{figure}{1}
\setcounter{secnumdepth}{4}


\setlength{\parskip}{1ex plus 0.5ex minus 0.2ex}


\usepackage{multirow}

\usepackage{setspace}
% \renewcommand{\baselinestretch}{10}
\newcommand{\mycaptiontable}[1]{
\begin{spacing}{0.6}
  %\vspace{0.5cm}
  \begin{quote}
    %\begin{center}
            {{Table} \thechapter.\arabic{table}: #1}
    %\end{center}
  \end{quote}
  %\vspace{1cm}
  \end{spacing}
  \stepcounter{table}
}

\newcommand{\mycaptionfigure}[1]{
  %\vspace{0.5cm}
  \begin{spacing}{0.6}
  \begin{quote}
    %\begin{center}
            {{Figure} \thechapter.\arabic{figure}: #1}
    %\end{center}
  \end{quote}
  %\vspace{1cm}
  \end{spacing}
  \stepcounter{figure}
}

% \newlength{\myVSpace}% the height of the box
% \setlength{\myVSpace}{0ex}% the default, 
% \newcommand\xstrut{\raisebox{-.5\myVSpace}% symmetric behaviour, 
%   {\rule{0pt}{\myVSpace}}%
% }

\usepackage{amsmath}

\usepackage{courier}

%***************************************************************************
%***************************************************************************
%***************************************************************************
\usepackage{multirow}
\usepackage{rotating}
\usepackage{setspace, amssymb, amsmath, epsfig, multirow, colortbl, tabularx}%
% For acronym package:
% If footnote is specified, text will be included in a footnote
% If printonlyused is specified, only used acronyms will be included
% I use the acronym sty under the sty directory as I needed the newest version
%\usepackage[footnote,printonlyused,withpage]{acronym} 
%\usepackage[printonlyused]{sty/acronym}

% Mejor que el acronym es glossaries
\usepackage[acronym,shortcuts,nomain]{glossaries}

% ifthen to allow using language dependent settings
\usepackage{ifthen}

\newcommand{\clearemptydoublepage}{\newpage{\pagestyle{empty}\cleardoublepage}}


\pagestyle{fancy}

% Pantone 160
%\definecolor{headingPortadaTFM}{RGB}{158,84,10}
% Pantone 160C (this is supposed to be the correct one, but it looks horrible in screen)
%\definecolor{headingPortadaTFM}{RGB}{161,86,28}
% Gold in RGB
%\definecolor{textoHeadingPortadaTFM}{RGB}{215,215,0}
% Captured colors in screen (this looks pst on screen)
\definecolor{headingPortadaTFM}{RGB}{152,118,52}
\definecolor{textoHeadingPortadaTFM}{RGB}{208,205,102}

\definecolor{azul}{rgb}{0,0,1}
\definecolor{verde}{rgb}{0,0.5,0}
\definecolor{rojo}{rgb}{1,0,0}
\definecolor{cyan}{rgb}{0,0.75,0.75}
\definecolor{magenta}{rgb}{0.75,0,0.75}
\definecolor{amarillo}{rgb}{0.75,0.75,0}
\definecolor{gris}{rgb}{0.25,0.25,0.25}
\definecolor{r}{rgb}{0,0,1}
\definecolor{g}{rgb}{0,1,0}
\definecolor{b}{rgb}{1,0,0}
\definecolor{c}{rgb}{0,1,1}
\definecolor{m}{rgb}{1,0,1}
\definecolor{y}{rgb}{1,1,0}
\definecolor{w}{rgb}{1,1,1}
\definecolor{k}{rgb}{0,0,0}
\definecolor{azulE}{rgb}{0,0.3984,0.5977}
\definecolor{amarilloE}{rgb}{0.9961,0.7969,0}
\definecolor{blanco}{rgb}{1,1,1}
\definecolor{burdeos}{rgb}{1,0,0.95}


\providecommand\phantomsection{}
\onehalfspacing
\sloppy  %better line breaks

\renewcommand{\chaptermark}[1]{\markboth{\chaptername\ \thechapter.\ #1}{}}
\renewcommand{\sectionmark}[1]{\markright{\thesection\ #1}{}}

%%%%%%%%%%%%%%%%%%%%%%%%%%%%%%%%%%%%%%%%%%%%%%%%%%%%%%%%%%%%%%%%%%%%%%%%%%%
% BEGIN Fancy headers stuff
\fancyhf{}

\fancyhead[LE,RO]{\bfseries\thepage}
\fancyhead[LO]{\bfseries\rightmark}
\fancyhead[RE]{\bfseries\leftmark}

\renewcommand{\headrulewidth}{0.5pt}
\renewcommand{\footrulewidth}{0pt}
\addtolength{\headheight}{3.5pt}
\fancypagestyle{plain}{\fancyhead{}\renewcommand{\headrulewidth}{0pt}}
\fancypagestyle{myplain}
{
  \fancyhf{}
  \renewcommand\headrulewidth{0pt}
  \renewcommand\footrulewidth{0pt}
  \fancyfoot[C]{\thepage}
}
% END Fancy headers stuff
%%%%%%%%%%%%%%%%%%%%%%%%%%%%%%%%%%%%%%%%%%%%%%%%%%%%%%%%%%%%%%%%%%%%%%%%%%%

%%%%%%%%%%%%%%%%%%%%%%%%%%%%%%%%%%%%%%%%%%%%%%%%%%%%%%%%%%%%%%%%%%%%%%%%%%%
% BEGIN Set nice chapter titles

% BEGIN Example 0 from http://texblog.org/2012/07/03/fancy-latex-chapter-styles/
% \usepackage[explicit]{titlesec}
% \usepackage{blindtext}
% \definecolor{gray75}{gray}{0.75}
% \newcommand{\hsp}{\hspace{20pt}}
% \titleformat{\chapter}[hang]{\Huge\bfseries}{\chaptername~\thechapter\hsp\textcolor{gray75}{|}\hsp}{0pt}{\Huge\bfseries}
% END Example 0 from http://texblog.org/2012/07/03/fancy-latex-chapter-styles/

% BEGIN Example 1 from http://texblog.org/2012/07/03/fancy-latex-chapter-styles/
% \usepackage{titlesec}
% \usepackage{blindtext}
% \definecolor{gray75}{gray}{0.75}
% \newcommand{\hsp}{\hspace{20pt}}
% \titleformat{\chapter}[hang]{\Huge\bfseries}{\chaptername~\thechapter\hsp\textcolor{gray75}{|}\hsp}{0pt}{\Huge\bfseries}
% END Example 1 from http://texblog.org/2012/07/03/fancy-latex-chapter-styles/

% BEGIN Example 2 from http://texblog.org/2012/07/03/fancy-latex-chapter-styles/
%Options: Sonny, Lenny, Glenn, Conny, Rejne, Bjarne, Bjornstrup
%\usepackage[Sonny]{fncychap}
%\usepackage[Lenny]{fncychap} % ugly
%\usepackage[Glenn]{fncychap}
%\usepackage[Conny]{fncychap} % ugly
%\usepackage[Rejne]{fncychap}
%\usepackage[Bjarne]{fncychap} % Doesn't work in Spanish
%\usepackage[Bjornstrup]{fncychap}
% END   Example 2 from http://texblog.org/2012/07/03/fancy-latex-chapter-styles/

% BEGIN Example 3 from http://texblog.org/2012/07/03/fancy-latex-chapter-styles/
% This is a nice colored example
% \usepackage{kpfonts}
% \usepackage[explicit]{titlesec}
% \newcommand*\chapterlabel{}
% \titleformat{\chapter}
%   {\gdef\chapterlabel{}
%    \normalfont\sffamily\Huge\bfseries\scshape}
%   {\gdef\chapterlabel{\thechapter\ }}{0pt}
%   {\begin{tikzpicture}[remember picture,overlay]
%     \node[yshift=-3cm] at (current page.north west)
%       {\begin{tikzpicture}[remember picture, overlay]
%         \draw[fill=LightSkyBlue] (0,0) rectangle
%           (\paperwidth,3cm);
%         \node[anchor=east,xshift=.9\paperwidth,rectangle,
%               rounded corners=20pt,inner sep=11pt,
%               fill=MidnightBlue]
%               {\color{white}\chapterlabel#1};
%        \end{tikzpicture}
%       };
%    \end{tikzpicture}
%   }
% \titlespacing*{\chapter}{0pt}{50pt}{-60pt}
% END   Example 3 from http://texblog.org/2012/07/03/fancy-latex-chapter-styles/

% BEGIN Example 4 from http://texblog.org/2012/07/03/fancy-latex-chapter-styles/
% END   Example 4 from http://texblog.org/2012/07/03/fancy-latex-chapter-styles/




% END Set nice chapter titles
%%%%%%%%%%%%%%%%%%%%%%%%%%%%%%%%%%%%%%%%%%%%%%%%%%%%%%%%%%%%%%%%%%%%%%%%%%%

%%%%%%%%%%%%%%%%%%%%%%%%%%%%%%%%%%%%%%%%%%%%%%%%%%%%%%%%%%%%%%%%%%%%%%%%%%%
% This is to set background images (in our case to set background image
% in TFMs front and back pages)
% If you want to set this background, use \BgThispage in the
% corresponding pages
\usepackage[pages=some]{sty/background}
\backgroundsetup{ scale=1, angle=0, opacity=.1, color=pink,
  contents={\includegraphics[width=.7\paperwidth]{logos/logoEPS-UAH.jpg}}, vshift=-50pt,  hshift=0pt }


% This is to allow do a clearpage and let the next one to be placed in
% even pages (to set a backpage for example)
\makeatletter
\newcommand*{\cleartoleftpage}{%
  \clearpage
    \if@twoside
    \ifodd\c@page
      \hbox{}\newpage
      \if@twocolumn
        \hbox{}\newpage
      \fi
    \fi
  \fi
}
\makeatother

